\documentclass[12pt]{article}

\input{0a - preambule}

\usepackage{pgfplotstable}
\usepackage{pgfmath}
\usepackage[utf8]{inputenc}
\usepackage[T1]{fontenc}
\usepackage{textcomp}

\title{\textit{Pourquoi} l'exponentielle}
\begin{document}

\maketitle 

\section{Jakob Bernoulli, 1654 - 1705}
Que devient 1 Thaler (monnaie à l'époque en Suisse) placé à 100\% d’intérêt par an, si on capitalise de plus en plus souvent ?

Si on capitalise :

\begin{itemize}
    \item une fois par an: $1 \text{ Thaler}\times (1 + 100\%) = (1 + 1) = 2  \text{ Thaler}$
    \item deux fois par an: $1 \text{ Thaler} \times (1 + 50\%)(1 + 50\%) = 1 \times (1 + 50\%)^2 = (1 + \frac{1}{2})^2 = 2,25\text{ Thaler} $
    \item trois fois par an: $1 \times (1 + \frac{100\%}{3})^3 = (1 + \frac{1}{3})^3$
    \item n fois par an: $1 \times (1 + \frac{100\%}{n})^n = (1 + \frac{1}{n})^n$
\end{itemize}

Et si on capitalise en continu?
\begin{itemize}
    \item 1 Thaler: $\lim_{n\to\infty}(1+\frac{1}{n})^n$
    \item $x$ Thaler: $\lim_{n\to\infty}(1+\frac{x}{n})^n$
\end{itemize}

Bernouilli a ainsi approximé la valeur pour $x=1$. 

On voit qu'elle converge rapidement:
\def\x{1}   % mets 1 pour voir e, 0 donnerait une erreur nulle

% --- Génération des données ---
\pgfplotstablenew[
    columns={n, val},
    create on use/n/.style={
        create col/expr={\pgfplotstablerow + 1}
    },
    create on use/val/.style={
        create col/expr={(1+\x/(\pgfplotstablerow+1))^(\pgfplotstablerow+1)}
    },
]{100}\bernoulidata

% --- Graphe montrant la convergence ---
\begin{center}
\begin{tikzpicture}
    \begin{axis}
        [
            width=12cm, height=7cm,
            xlabel={$n$},
            ylabel={$\left(1+\frac{\x}{n}\right)^n$},
            xmin=1, xmax=100,
            grid=both,
        ]
        \addplot+[mark=*, mark size=1.2pt] table[x=n, y=val]{\bernoulidata};
    \end{axis}
\end{tikzpicture}
\end{center}


\begin{gather*}
    \lim_{n\to\infty}(1+\frac{1}{n})^n = \approx 2,718281828459045
\end{gather*}

Cette fonction a été nommée par la suite exponentielle:
\begin{gather*}
    e^x=\lim_{n\to\infty}(1+\frac{x}{n})^n
\end{gather*}

\paragraph{On remarque déjà que $(e^x)^{(n)}=e^x$}
Soit
\[
f(x)=\left(1+\frac{x}{n}\right)^n
\]

On pose
\[
u(x)=1+\frac{x}{n}
\]
Alors
\[
f(x)=u(x)^n
\]

Par la règle de la chaîne,
\[
f'(x)=n\,u(x)^{\,n-1}\,u'(x).
\]

Or
\[
u'(x)=\frac{1}{n}
\]

Donc
\[
f'(x)=n\left(1+\frac{x}{n}\right)^{n-1}\cdot\frac{1}{n}
=\left(1+\frac{x}{n}\right)^{n-1}
\]

Ainsi,
\[
\boxed{
\frac{d}{dx}\left(1+\frac{x}{n}\right)^n
=\left(1+\frac{x}{n}\right)^{n-1}
}
\]

donc, lorsque n tend vers l'infini:
\[
    (e^x)'=e^x
\]


\section{Leonhard Euler 1707 - 1783}
Il connaissait les séries de Taylor (Brooke Taylor 1685 - 1731) pour approximer les fonctions suffisamment dérivables autour d'un voisinage $a$:
\begin{gather*}
    f(x)=\sum_{n=0}^{\infty}\frac{f^{(n)}(a)}{n!}(x-a)^n
\end{gather*}

Comme $f^{(n)}(a)$ et $n!$ sont des constantes, on dérive seulement $(x-a)^n$

\[
f'(x)
=\sum_{n=1}^{\infty}\frac{f^{(n)}(a)}{n!}\,n(x-a)^{n-1}
\]

\[
f'(x)=\sum_{n=1}^{\infty}\frac{f^{(n)}(a)}{(n-1)!}(x-a)^{n-1}
\]

\[
\text{En posant } k=n-1,\ \ f'(x)=\sum_{k=0}^{\infty}\frac{f^{(k+1)}(a)}{k!}(x-a)^k = f(x) \text{ si } f^{(k+1)} = f^{(k)}
\]


L'exponentielle est une fonction tellement régulière (son rayon de convergence est infini) que la formule de Taylor fonctionnera parfaitement quel que soit $a$.

Lorsque $a$ = 0:
\begin{gather*}
    f(x)=\sum_{n=0}^{\infty}\frac{f^{(n)}(0)}{n!}x^n
\end{gather*}

Dans le cas de $e^x$, toutes ses dérivées sont égales: $e^x=(e^x)'=(e^x)''=...=(e^x)^{(n)}$

Et pour $x=0$, elles valent toutes $1$. Autrement dit:
\begin{gather*}
    \forall n \in \mathbb{N^*}, \quad f^{(n)}(0)=1
\end{gather*}

D'où
\begin{gather*}
    e^x=\sum_{n=0}^{\infty}\frac{f^{(n)}(0)}{n!}x^n=\sum_{n=0}^{\infty}\frac{1}{n!}x^n \\
    \boxed{e^x=\sum_{n=0}^\infty{\frac{x^n}{n!}}}
\end{gather*}

Illustration de la convergence de la série:
% La fonction sum n'est pas disponible
\pgfmathdeclarefunction{s_1}{1}{\pgfmathparse{1 + #1}}
\pgfmathdeclarefunction{s_2}{1}{\pgfmathparse{1 + #1 + (#1^2)/2}}
\pgfmathdeclarefunction{s_3}{1}{\pgfmathparse{1 + #1 + (#1^2)/2 + (#1^3)/6}}
\pgfmathdeclarefunction{s_5}{1}{%
  \pgfmathparse{1 + #1 + (#1^2)/2 + (#1^3)/6 + (#1^4)/24 + (#1^5)/120}%
}
\pgfmathdeclarefunction{s_10}{1}{%
  \pgfmathparse{1 + #1 + (#1^2)/2 + (#1^3)/6 + (#1^4)/24 + (#1^5)/120
    + (#1^6)/720 + (#1^7)/5040 + (#1^8)/40320 + (#1^9)/362880 + (#1^10)/3628800}%
}

\begin{center}
\begin{tikzpicture}
    \begin{axis}[
            axis lines=middle,
            width=10cm,
            height=10cm,
            xlabel={$x$}, ylabel={$y$},
            grid=both,
            xmin=-6, xmax=5,
            ymin=-5, ymax=30,
            samples=25,
            domain=-6:5,
            legend style={at={(0.02,0.98)}, anchor=north west, draw=none},
            ]
           
        \addplot[mark=x, draw=gray] {s_1(x)};
        \addlegendentry{$n=1$}
        
        \addplot[mark=+, draw=brown] {s_2(x)};
        \addlegendentry{$n=2$}
        
        \addplot[mark=x, draw=orange] {s_5(x)};
        \addlegendentry{$n=5$}
        
        \addplot[mark=+, draw=red] {s_10(x)};
        \addlegendentry{$n=10$}

        \addplot[mark=x, draw=blue] {exp(x)};
        \addlegendentry{$e^x$}
        
    \end{axis}
\end{tikzpicture}
\end{center}


Vérifions le. \\
Pour rappel, $n!$, qui se prononce \textit{factorielle n}, est une fonction définie pour tout $n\in \mathbb{N}$ ainsi:

\begin{gather*}
    0!=1 \\
    n!=1\times2\times3\times4\times5\times...(n-1)\times n
\end{gather*}

On peut aussi la définir récursivement par 
\begin{gather*}
    0!=1 \\
    n!=n(n-1)!
\end{gather*}
ce qui va nous être utile par la suite.

Prouvons que \textbf{DOUBLON? ou version 2 pour la pédagogie?}
\begin{gather*}
    e^x=f(x)=\sum_{n=0}^\infty{\frac{x^n}{n!}}
\end{gather*}

Pour rappel, $0^0 = 1$, car on ne multiplie aucune fois par $0$. On a donc:
\begin{gather*}
    f(0)=\sum_{n=0}^\infty{\frac{0^n}{n!}}=\frac{0^0}{0!}+\frac{0^1}{1!}+\frac{0^2}{2!}+... \\
    \boxed{$f(0)=1$}
\end{gather*}


Prouvons maintenant que $f(x)=f'(x)$
\begin{gather*}
    f(x)=\sum_{n=0}^\infty{\frac{x^n}{n!}} \\
    f(x)=\frac{x^0}{0!}+\frac{x^1}{1!}+\frac{x^2}{2!}+\frac{x^3}{3!}+...+\frac{x^n}{n!}+... \\
\end{gather*}

Développons $f(x)$ pour mieux comprendre ce qui va se passer lorsqu'on va dériver $f$:
\begin{gather*}
    f(x)=1+x+\frac{x^2}{1 \times2}+\frac{x^3}{1 \times2 \times3}+...+\frac{x^n}{1 \times2  \times3  \times ... \times n} \\
    f'(x)=0+1+2\frac{x^{2-1}}{2!}+3\frac{x^{3-1}}{3!}+...+n\frac{x^{n-1}}{n!}+... \\
    f'(x)=1+\frac{x^{1}}{1!}+\frac{x^{2}}{2!}+...+\frac{x^{n-1}}{(n-1)!}+...=f(x)
\end{gather*}

Lorsqu'on dérive $f$, on retombe sur la même fonction, tous les éléments sont simplement décalés d'un cran à gauche, mais puisqu'il y en a une infinité, rien n'a changé:

Donc $\boxed{f'(x)=f(x)}$

\end{document}
